% Figure shorelines: generated by figures/rse_fig_shorelines.py; placed where the text first discusses it.
% Choice of the level bin, waterline cells, definition of Q, graticule and scale bars: Supplementary
% section rse_si_display.tex.
\begin{figure}[!tbp]
  \centering
  \includegraphics[width=\textwidth,height=0.66\textheight,keepaspectratio]{../figures/out/rse_fig_shorelines.pdf}
  \caption{\textbf{Where the five grids put the shoreline.}
  All panels show the level bin centred at \metric{erl.hero.e3_b029.level_m}~m (DAHITI reference), observed by \metric{erl.hero.e3_b029.n_acquisitions} Landsat acquisitions in \metric{erl.hero.e3_b029.first_year}--\metric{erl.hero.e3_b029.last_year}.
  (a)~Observed sea (blue), land (beige) and cells of undetermined state (white), the observed shoreline cells (black), the zooms of c--e and the transect of b (vermilion).
  (b)~Profile along the transect on the Volga--Ural interfluve coast, which follows one ICESat-2 ATL08 beam from sea to land.
  Coloured lines: the five grids' heights after each grid's area-fitted vertical offset; black line: the waterline elevation model on the emerged seabed; grey dots: ATL08 heights; dashed line: the level.
  Ticks under the profile: where the observed shoreline (black) and each grid's predicted shoreline cross the transect; dotted vertical lines repeat the observed crossings.
  (c--e)~Zooms on the three checked coasts, one panel per grid: observed shoreline cells (black) and every predicted shoreline cell, coloured by its distance to the nearest observed shoreline cell; grey, predicted shoreline cells of unknown observed state, left out of the displacement.
  In each panel, $D$ is the grid's median displacement in the sector at this level (both directions pooled); $Q$, in the row heading, is the sector's observation noise floor at this level (Fig.~\ref{fig:rse_method}c).}
  \label{fig:rse_shorelines}
\end{figure}
